% GENERATED FILE. Edit canonical lessons, then run npm run generate:books.
% canonical-source-hash: fnv1a64:87f7656e83a12e31
% canonical-lessons: LA-C283-puto, LA-C283-mea-sententia, LA-C283-ergo, LA-R283-what-i-think

\chapter{What I Think}
\label{ch:la-what-i-think}
\hlchaptermodality{283}

\begin{chapteropening}
\textbf{By the end of this chapter:} \emph{I can say what I think, with putō and an accusative and infinitive, and draw a conclusion with ergō.}

Complete LA-R283-what-i-think: I can say what I think, with putō and an accusative and infinitive, and draw a conclusion with ergō.
\end{chapteropening}

\section[putō, putāre]{putō — I think that …}

\label{lesson:LA-C283-puto}



\begin{quote}
Say \textquotedblleft{}I think\textquotedblright{}, the verb of turning things over. (\emph{Cōgitō.}) Now \textquotedblleft{}I believe\textquotedblright{}. (\emph{Crēdō.})

The chapter on \emph{cōgitō} named a lighter verb beside it, closer to \textquotedblleft{}I reckon\textquotedblright{}. This lesson puts it to work.
\end{quote}

\subsection*{You'll want to know: putō}
\textbf{putō, putāre} — \textquotedblleft{}to think, to reckon, to suppose\textquotedblright{}. It is the everyday verb of opinion: \emph{cōgitō} is thinking hard; \emph{putō} is what you think.

\begin{quote}
\textbf{Ita putō.} — I think so; that is how I reckon it.
\end{quote}

Its perfect is regular: \textbf{putāvī}, I thought.

\begin{grammarlens}[title={Grammar Lens: no word for \textquotedblleft{}that\textquotedblright{}}]
English says \textquotedblleft{}I think \textbf{that} my friend is ill\textquotedblright{}. Latin has no word there. It puts the person in the \textbf{accusative} and the verb in the \textbf{infinitive}:

\begin{quote}
\textbf{Putō amīcum aegrum esse.} — I think my friend is ill.
\end{quote}

Word for word, \textquotedblleft{}I think friend-\emph{um} to be ill\textquotedblright{}. The \emph{-um} of \emph{amīcum} marks who the thought is about, and \emph{esse} carries the verb. Even the adjective joins the accusative: \emph{aegrum}, not \emph{aeger}.

\begin{quote}
\textbf{Crēdō tē fessum esse.} — I believe you are tired. \textbf{Putō pluere.} — I think it is raining.
\end{quote}

\emph{Pluere} has no subject at all, as \emph{pluit} had none. This pattern, accusative and infinitive, carries what someone thinks, believes or says all through Latin.
\end{grammarlens}

\begin{cousinweb}
The \emph{cōgitō} lesson told the story: \textbf{putō} first meant \textquotedblleft{}to prune\textquotedblright{}, then \textquotedblleft{}to settle an account\textquotedblright{}, then \textquotedblleft{}to think\textquotedblright{} (Lewis \& Short trace it to a root meaning \textquotedblleft{}to clean\textquotedblright{}). You own two of its compounds already: \emph{computō}, \textquotedblleft{}to calculate\textquotedblright{}, is \emph{com-} + \emph{putō}, reckoning together, and \emph{disputō}, \textquotedblleft{}to discuss\textquotedblright{}, began in the merchants' sense of going over a sum item by item.

The Roman writer Aulus Gellius made a point of it. He read the word for \textquotedblleft{}I think\textquotedblright{} as a pruner's word: cutting away the false opinions until what is left is true (Lewis \& Short quote the passage under \emph{puto}).
\end{cousinweb}

\subsection*{Your turn}
Say these aloud:

\begin{itemize}
\raggedright
  \item \emph{putō} — then \emph{Ita putō.}
  \item \emph{Putō amīcum aegrum esse.}
  \item \emph{Crēdō tē fessum esse.}
  \item \emph{Putō pluere.}
  \item \emph{Herī putāvī amīcum venīre.} — Yesterday I thought my friend was coming
\end{itemize}

\begin{tcolorbox}[breakable,colback=teal!4,colframe=teal!35!black,title={Before you move on}]
How does Latin say \textquotedblleft{}I think my friend is ill\textquotedblright{}? (\emph{\textbf{Putō amīcum aegrum esse.}}) Which case does the person go into? (The \textbf{accusative}.) And the verb? (The \textbf{infinitive}.) What did \emph{putō} first mean? (\textbf{To prune.})
\end{tcolorbox}
\section[meā sententiā, mihi vidētur]{meā sententiā, mihi vidētur — in my opinion}

\label{lesson:LA-C283-mea-sententia}



\begin{quote}
Say \textquotedblleft{}I think my friend is ill\textquotedblright{}. (\emph{Putō amīcum aegrum esse.})

\emph{Putō} turns the whole sentence into what you think. Latin also has ways to leave the sentence as it is and hang the opinion on it.
\end{quote}

\subsection*{You'll want to know: meā sententiā}
\textbf{meā sententiā} — \textquotedblleft{}in my opinion\textquotedblright{}. You own \textbf{sententia}, \textquotedblleft{}an opinion\textquotedblright{}; this is its ablative with \emph{meā}, \textquotedblleft{}my\textquotedblright{}: word for word, \textquotedblleft{}by my opinion\textquotedblright{}.

\begin{quote}
\textbf{Meā sententiā, Rōma pulchra est.} — In my opinion, Rome is beautiful.
\end{quote}

Nothing else in the sentence changes: \emph{Rōma} stays nominative, \emph{est} stays as it was. Plautus says it, and so does Cicero.

\subsection*{You'll want to know: mihi vidētur}
\textbf{mihi vidētur} — \textquotedblleft{}it seems to me\textquotedblright{}. You own \textbf{videor}, \textquotedblleft{}to seem\textquotedblright{}; \emph{mihi} is \textquotedblleft{}to me\textquotedblright{}.

\begin{quote}
\textbf{Mārcus mihi fessus vidētur.} — Marcus seems tired to me. \textbf{Pluit, ut mihi vidētur.} — It is raining, as far as I can tell.
\end{quote}

Marcus is the subject, so \emph{fessus} stays nominative. And \emph{ut mihi vidētur}, \textquotedblleft{}as it seems to me\textquotedblright{}, is a phrase Cicero tucks into the middle of a sentence.

\begin{cousinweb}
\textbf{sententia} is built on \emph{sentiō}, \textquotedblleft{}to feel, to think\textquotedblright{} (Lewis \& Short), so \emph{meā sententiā} is \textquotedblleft{}by what I feel about it\textquotedblright{}. \textbf{vidētur} is the passive of \emph{videō}, \textquotedblleft{}to see\textquotedblright{}: \textquotedblleft{}it is seen\textquotedblright{}. So \emph{mihi vidētur} is, word for word, \textquotedblleft{}it is seen, to me\textquotedblright{}, which is how a Roman said \textquotedblleft{}it seems to me\textquotedblright{}. \emph{Videō} is inherited from Proto-Indo-European \emph{*weyd-}, \textquotedblleft{}to see\textquotedblright{} (EtymDB 2.1).
\end{cousinweb}

\subsection*{Your turn}
Say these aloud:

\begin{itemize}
\raggedright
  \item \emph{meā sententiā}, then \emph{mihi vidētur}
  \item \emph{Meā sententiā, Rōma pulchra est.}
  \item \emph{Mārcus mihi fessus vidētur.}
  \item the same opinion three ways — \emph{Putō Rōmam pulchram esse. Meā sententiā, Rōma pulchra est. Rōma mihi pulchra vidētur.}
\end{itemize}

\begin{tcolorbox}[breakable,colback=teal!4,colframe=teal!35!black,title={Before you move on}]
\textquotedblleft{}In my opinion\textquotedblright{}? (\emph{\textbf{Meā sententiā.}}) \textquotedblleft{}It seems to me\textquotedblright{}? (\emph{\textbf{Mihi vidētur.}}) After \emph{putō}, \emph{Rōma} becomes what? (\emph{\textbf{Rōmam}}, with \emph{esse}.) After \emph{meā sententiā}? (It stays \emph{\textbf{Rōma}}.)
\end{tcolorbox}
\section[ergō, itaque, igitur]{ergō, itaque, igitur — and so}

\label{lesson:LA-C283-ergo}



\begin{quote}
The \emph{cōgitō} lesson gave you a famous sentence. Say it. (\emph{Cōgitō, ergō sum.} I think, therefore I am.) The sentence is Descartes', written long after Rome, but every word of it is Latin.

You also own \textbf{itaque}, \textquotedblleft{}and so\textquotedblright{}, and \textbf{igitur}, \textquotedblleft{}therefore\textquotedblright{}. This lesson sorts all three.
\end{quote}

\begin{grammarlens}[title={Grammar Lens: what happened next, and what follows}]
\textbf{itaque} tells a story forward. It stands first, and what comes after it is the next thing that happened, because of the last:

\begin{quote}
\textbf{Pluēbat; itaque in vīllā mānsimus.} — It was raining, and so we stayed in the farmhouse.
\end{quote}

\textbf{igitur} and \textbf{ergō} draw a conclusion. They belong to an argument more than to a story. In Cicero's prose \emph{igitur} usually comes second, after the first word: \emph{Fessus igitur es.} — So you are tired.

\textbf{ergō}, \textquotedblleft{}therefore, so\textquotedblright{}, often stands first, and it has the sound of someone reaching the point: \emph{Fessus es; ergō dormī.} — You are tired; so sleep.

So: a story's \emph{and so} is \emph{itaque}; an argument's \emph{therefore} is \emph{igitur} or \emph{ergō}.
\end{grammarlens}

\begin{cousinweb}
Before \textbf{ergō} meant \textquotedblleft{}therefore\textquotedblright{}, it meant \textquotedblleft{}on account of\textquotedblright{}, and it followed a genitive: \emph{virtūtis ergō}, \textquotedblleft{}on account of courage\textquotedblright{}, an old formula for decorating a soldier (Lewis \& Short). Lewis \& Short build the word from \emph{ex}, \textquotedblleft{}out of\textquotedblright{}, and a root meaning \textquotedblleft{}to stretch, to direct\textquotedblright{}, but the origin is not settled. The sense moved from \textquotedblleft{}on account of this\textquotedblright{} to \textquotedblleft{}so, then\textquotedblright{}.
\end{cousinweb}

\subsection*{Your turn}
Say these aloud:

\begin{itemize}
\raggedright
  \item \emph{Cōgitō, ergō sum.}
  \item a story — \emph{Pluēbat; itaque in vīllā mānsimus.}
  \item an argument — \emph{Fessus igitur es.} — then \emph{Fessus es; ergō dormī.}
  \item \emph{Meā sententiā, amīcus aeger est; ergō nōn venit.}
\end{itemize}

\begin{tcolorbox}[breakable,colback=teal!4,colframe=teal!35!black,title={Before you move on}]
Which \textquotedblleft{}so\textquotedblright{} carries a story forward? (\emph{\textbf{Itaque.}}) Which two draw a conclusion? (\emph{\textbf{Igitur}} and \emph{\textbf{ergō}}.) Where does \emph{igitur} usually stand in Cicero? (\textbf{Second}, after the first word.) And what did \emph{ergō} mean before \textquotedblleft{}therefore\textquotedblright{}? (\textbf{On account of}, after a genitive.)
\end{tcolorbox}
\section[(review)]{Review — what I think, and what follows}

\label{lesson:LA-R283-what-i-think}



\begin{quote}
Say \textquotedblleft{}I think it is raining\textquotedblright{}. (\emph{Putō pluere.}) Then \textquotedblleft{}I believe you are tired\textquotedblright{}. (\emph{Crēdō tē fessum esse.})
\end{quote}

\begin{grammarlens}[title={What you've built}]
Three lessons, and an opinion has three shapes. \textbf{putō} with the accusative and infinitive, which turns the sentence into what you think. \textbf{meā sententiā}, which leaves it alone. And \textbf{mihi vidētur}, \textquotedblleft{}it seems to me\textquotedblright{}, with the thing that seems as its subject. Then \textbf{ergō} to draw the conclusion, beside the \emph{itaque} of a story and the \emph{igitur} of an argument.

\begin{quote}
\emph{Meā sententiā, Mārcus fessus est; ergō dormit.} — In my opinion Marcus is tired; so he is sleeping.
\end{quote}

The opinion leaves \emph{Mārcus} nominative. Say it with \emph{putō} instead, and he turns accusative: \emph{Putō Mārcum fessum esse.}
\end{grammarlens}

\subsection*{Your turn}
Say these aloud:

\begin{itemize}
\raggedright
  \item one opinion three ways — \emph{Putō Rōmam pulchram esse. Meā sententiā, Rōma pulchra est. Rōma mihi pulchra vidētur.}
  \item \emph{Meā sententiā, Mārcus fessus est; ergō dormit.}
  \item a story's \textquotedblleft{}and so\textquotedblright{} — \emph{Pluēbat; itaque in vīllā mānsimus.}
  \item an opinion of your own about today's weather, with \emph{putō}, then a conclusion with \emph{ergō}
\end{itemize}

\begin{tcolorbox}[breakable,colback=teal!4,colframe=teal!35!black,title={Before you move on}]
\textquotedblleft{}In my opinion\textquotedblright{}? (\emph{\textbf{Meā sententiā.}}) \textquotedblleft{}It seems to me\textquotedblright{}? (\emph{\textbf{Mihi vidētur.}}) What did \emph{putō} first mean? (\textbf{To prune.}) After \emph{putō}, which case is the person in? (The \textbf{accusative}.) And from the last chapter: what does \emph{dum} mean? (\textbf{While.}) Which rung comes first? (\emph{\textbf{Prīmum.}})
\end{tcolorbox}
